%% =====================================================================
%%  B3-T-08-01 -- what B3 contributes to "Do Not Offend the Wrong Person"
%%  -------------------------------------------------------------------
%%  LAYER A: APPEND-ONLY. Written once, then left alone. Material found
%%  only here sits in the `unique` slot and is protected by rule R10.
%% =====================================================================
%% @kind: source
%% @id: B3-T-08-01
%% @book: B3
%% @topic: T-08-01
%% @locator: Law 19, pp. 137--147
%% @status: distilled
%% @unique: yes
%% @integrated: yes
%% @updated: 2026-09-19

\dossier{B3}{T-08-01}{\bookshort{B3}}{Law 19, pp. 137--147}

\begin{distilled}
  \item There are many kinds of people and you cannot assume everyone reacts to your strategies the same way.
  \item Some people, once deceived or outmanoeuvred, spend the rest of their lives seeking revenge; they are described as wolves in lambs' clothing.
  \item The instruction is to choose victims and opponents carefully, and never to offend or deceive the wrong person.
  \item The law is supported by a long case of a ruler who insulted a rising neighbour repeatedly over a small matter and lost an empire as a result.
\end{distilled}

\begin{unique}
  \item The delayed-revenge disposition named as the primary strategic risk, with the consequence that the offence and the cost are separated in time and the offender does not connect them.
  \item The wolves-in-lambs'-clothing formulation: the dangerous counterpart presents as harmless at the time of the offence.
\end{unique}
